% mazzi: 29
\chapter{Il sistema nervoso periferico}

% ---------------------------------------------------------------------------
\section{Generalità}

\textbf{Funzione}: il SNP permette la connessione di encefalo e midollo spinale
(\textit{nevrasse}) con la periferia.

\textbf{Componenti principali}: nervi spinali, nervi cranici, gangli, organi di
senso.

\textbf{Nervi somatici}: emergono dal nevrasse e si distribuiscono ai distretti
superficiali e scheletrici del corpo (cute, alcune mucose, ossa, muscoli,
articolazioni e organi di senso).

I nervi sono distinti in sensitivi e motori in base alla tipologia di stimolo
che trasportano; si individuano 4 categorie:
\begin{itemize}
  \item nervi sensitivi somatici;
  \item nervi sensitivi viscerali;
  \item nervi motori somatici;
  \item nervi effettori (motori) viscerali.
\end{itemize}

\textbf{Nervi spinali}: presentano sempre tutte e quattro le componenti
(nervi misti).

\textbf{Nervi cranici}: possono presentarne una (nervi puri), due (nervi
semipuri) o quattro (nervi misti).

\rubrica{Collegamento verso la periferia}
\begin{itemize}
  \item \textbf{fibre sensitive}: portano informazioni dalla periferia al
    centro (\textit{fibre afferenti});
  \item \textbf{fibre effettrici o motorie}: dirette ai muscoli scheletrici,
    al miocardio, alla muscolatura liscia e alle ghiandole esocrine ed
    endocrine (\textit{fibre efferenti}).
\end{itemize}
Entrambi i tipi di fibre possono riferirsi sia a funzioni somatiche sia a
funzioni vegetative.

\rubrica{Nervi sensitivi e motori}
\begin{itemize}
  \item \textbf{Nervi sensitivi somatici}: raccolgono stimoli dai recettori
    periferici della sensibilità generale (pressione, temperatura, dolore) e
    della sensibilità specifica (5 sensi);
  \item \textbf{Nervi sensitivi viscerali}: trasportano impulsi provenienti
    dai diversi organi o sistemi;
  \item entrambi conducono gli impulsi al nevrasse.
\end{itemize}

\textbf{Nervi motori}: conducono impulsi in direzione centrifuga terminando
sulla muscolatura striata o sulla muscolatura liscia e sulle ghiandole. Le
\textbf{fibre effettrici viscerali} sono distinte in:
\begin{itemize}
  \item fibre motorie viscerali;
  \item fibre eccito-secretrici, con innervazione nei muscoli lisci e nelle
    ghiandole.
\end{itemize}

\rubrica{Gangli}
\begin{itemize}
  \item \textbf{Gangli sensitivi}: costituiti da neuroni sensitivi (ricezione
    di afferenze sensitive); situati sul decorso dei nervi cranici e delle
    radici posteriori dei nervi spinali;
  \item \textbf{Gangli autonomici}: costituiti da neuroni effettori
    viscerali; distinti in simpatici e parasimpatici.
\end{itemize}
Sono avvolti da una capsula connettivale; al loro interno si trovano fibre
nervose e cellule distinte in \textit{satelliti} (di natura gliale) e
\textit{principali}.

% ---------------------------------------------------------------------------
\section{Il nervo periferico}

I nervi spinali sono assimilabili a cordoni cilindrici biancastri,
costituiti da fasci paralleli di fibre nervose avvolti da guaine di tessuto
connettivo fibroso, con funzione protettiva e trofica. Procedendo dalla
superficie in profondità si distinguono:
\begin{itemize}
  \item \textbf{endonevrio}: sottile strato di tessuto connettivo lasso che
    circonda ogni singola fibra assonica con la sua cellula di Schwann di
    rivestimento;
  \item \textbf{perinevrio}: strato più spesso di tessuto connettivo, circonda
    ogni fascio;
  \item \textbf{epinevrio}: circonda i nervi periferici formati da più fasci,
    riempiendo anche lo spazio tra i vari fasci.
\end{itemize}

% ---------------------------------------------------------------------------
\section{I nervi spinali}

\textbf{Nervi spinali}: emergono in serie regolari dalla superficie laterale
del midollo spinale, attraverso l'unione di una radice posteriore
(sensitiva) e una anteriore (motoria), impegnandosi singolarmente nei forami
intervertebrali. Dalla confluenza delle due radici anteriore e posteriore si
forma il nervo spinale, che risulta pertanto misto.

Sono $33$ paia, suddivisi in $8$ cervicali, $12$ toracici, $5$ lombari, $5$
sacrali, $3$ coccigei; sono tutti nervi misti. Ai $33$ segmenti spinali
corrisponde una serie di segmenti cutanei (\textbf{dermatomeri}) e
muscolari (\textbf{miomeri}).

\rubrica{Dermatomeri e miomeri}
\textbf{Dermatomeri e miomeri}: aree e territori di innervazione periferica
che mostrano un regolare ordinamento metamerico o segmentale: ai $33$
segmenti spinali, ognuno con il proprio paio di nervi, corrispondono una
serie di segmenti cutanei e muscolari.

Ogni paio di nervi spinali controlla una specifica regione della superficie
cutanea. I dermatomeri sono clinicamente importanti, perché un danno a un
nervo spinale o a un ganglio della radice dorsale produce una caratteristica
perdita della sensibilità in uno specifico distretto cutaneo.

% ---------------------------------------------------------------------------
\section{I plessi nervosi}

\textbf{Plessi nervosi}: formati dai rami anteriori (ventrali) dei nervi
spinali che si anastomizzano tra loro; si distinguono in direzione
cranio-caudale.

\rubrica{Plesso cervicale}
Formato dai rami ventrali dei nervi spinali C1--C4 e da alcune fibre nervose
provenienti da C5. Da esso originano i nervi che innervano aree della testa,
del collo e del torace. Si individuano rami cutanei (sensitivi), rami
muscolari (motori) e rami anastomotici.

\rubrica{Plesso brachiale}
Formato dai rami ventrali dei nervi spinali C5--T1 (i rami anteriori dei
nervi cervicali da C5 a C8 e T1, insieme al ramo discendente di C4 e a una
minima componente ascendente di T2). I suoi rami innervano il cingolo
scapolare e l'arto superiore.

È localizzato nella regione sopraclavicolare tra i muscoli scaleno
anteriore, medio e posteriore, e presenta numerose anastomosi e
suddivisioni.
\begin{enumerate}
  \item \textbf{3 tronchi principali}:
  \begin{itemize}
    \item superiore (C4--C6);
    \item medio (C7);
    \item inferiore (C8--T2).
  \end{itemize}
  \item da questi derivano 2 rami:
  \begin{itemize}
    \item in prossimità del cavo ascellare, i rami posteriori di tutti i
      tronchi principali si uniscono a formare la \textbf{corda posteriore};
    \item i rami anteriori dei tronchi principali superiore e medio si
      fondono a costituire la \textbf{corda laterale};
    \item il ramo anteriore del tronco principale inferiore mantiene la
      propria individualità, costituendo la \textbf{corda mediale}.
  \end{itemize}
\end{enumerate}

Dalle tre corde derivano i rami terminali (sensitivi e motori), nonché rami
anastomotici e collaterali. I \textbf{rami terminali} sono $6$ e derivano
dai tronchi secondari:
\begin{itemize}
  \item nervo muscolo-cutaneo;
  \item nervo mediano;
  \item nervo ulnare;
  \item nervo cutaneo mediale dell'avambraccio;
  \item nervo cutaneo mediale del braccio;
  \item (oltre a rami anastomotici e rami collaterali anteriori e
    posteriori).
\end{itemize}

\rubrica{Plesso lombare}
Formato dai rami ventrali dei nervi spinali T12--L4. È situato profondamente
tra i piani di inserzione della porzione paravertebrale del muscolo psoas;
costituito da un complesso anastomotico a cui partecipano i rami anteriori
del $1^\circ$--$4^\circ$ nervo lombare, insieme al ramo anastomotico del
$12^\circ$ nervo intercostale per il $1^\circ$ nervo lombare. Ogni radice del
plesso si suddivide in 3 rami.

\rubrica{Plesso sacrale}
Formato dai rami ventrali dei nervi spinali L4--S4. È compreso tra la fascia
pelvica e il muscolo piriforme, all'interno della piccola pelvi, ed è
costituito dall'unione del tronco lombosacrale (L4, L5) e dai rami
anteriori del $1^\circ$, $2^\circ$ e $3^\circ$ nervo sacrale. Si distribuisce,
tramite i suoi numerosi rami (tra cui i nervi tibiale e peroniero comune),
alla muscolatura del compartimento posteriore della coscia, della gamba e
del piede, oltre che alla cute della gamba e del piede.

% ---------------------------------------------------------------------------
\section{I riflessi}

\textbf{Riflesso}: risposta automatica, involontaria, a uno stimolo
sensoriale. La via riflessa più semplice è detta \textbf{arco riflesso} e
consta di 5 componenti:
\begin{enumerate}
  \item un recettore sensoriale;
  \item un neurone afferente;
  \item un centro di integrazione (che può realizzarsi a livello del midollo
    spinale o a livello encefalico);
  \item un neurone efferente;
  \item un organo effettore.
\end{enumerate}

I riflessi possono essere:
\begin{itemize}
  \item \textbf{monosinaptici}: la via nervosa consta di 2 neuroni e 1
    sinapsi;
  \item \textbf{polisinaptici}: la via nervosa consta di più di 2 neuroni e
    più di 1 sinapsi.
\end{itemize}

\rubrica{Classificazione dei riflessi}
I riflessi si classificano inoltre:
\begin{itemize}
  \item in base allo sviluppo: \textbf{innati} (geneticamente determinati) o
    \textbf{acquisiti} (appresi);
  \item in base al sito di elaborazione: \textbf{spinali} (elaborati nel
    midollo spinale) o \textbf{cranici} (elaborati nell'encefalo);
  \item in base alla risposta motoria: \textbf{somatici} (controllano la
    contrazione della muscolatura scheletrica, includono i riflessi
    superficiali e da stiramento) o \textbf{viscerali/autonomi}
    (controllano l'azione della muscolatura liscia e cardiaca e delle
    ghiandole);
  \item in base alla complessità del circuito: \textbf{monosinaptici} (una
    sinapsi) o \textbf{polisinaptici} (da due a diverse centinaia di
    sinapsi).
\end{itemize}

\rubrica{Riflesso miotattico (o rotuleo)}
\begin{enumerate}
  \item lo stiramento del tendine stira il recettore sensoriale (fuso
    muscolare) del muscolo estensore della gamba;
  \item il neurone sensoriale eccita il motoneurone del midollo spinale e
    anche un interneurone spinale, che a sua volta inibisce il muscolo
    flessore;
  \item il motoneurone conduce il potenziale d'azione alla sinapsi delle
    fibre del muscolo estensore, provocando la contrazione; il muscolo
    flessore si rilascia perché l'attività dei suoi motoneuroni è stata
    inibita;
  \item la gamba si estende.
\end{enumerate}

\rubrica{Riflesso da stiramento}
Regolazione automatica della lunghezza del muscolo: lo stiramento del
muscolo stimola i fusi neuromuscolari, che attivano un neurone sensitivo;
l'informazione viene elaborata dal motoneurone, che si attiva e produce la
contrazione muscolare di risposta.

% ---------------------------------------------------------------------------
\section{I nervi cranici}

\textbf{Nervi cranici}: gruppo di fasci nervosi che originano direttamente
dal tronco encefalico; fanno parte del SNP; sono $12$ paia.

Forniscono innervazione motoria e sensoriale principalmente alle strutture
all'interno del cranio e del collo. L'innervazione sensoriale include sia
la sensazione ``generale'' (temperatura, tatto), sia l'innervazione
``specifica'' (gusto, vista, odore, equilibrio, udito).

\rubrica{I nervo -- Olfattivo}
\begin{itemize}
  \item \textbf{Funzione}: sensibilità specifica (olfatto);
  \item \textbf{origine}: recettori dell'epitelio olfattivo;
  \item \textbf{destinazione}: bulbi olfattivi.
\end{itemize}

\rubrica{II nervo -- Ottico}
\begin{itemize}
  \item \textbf{Funzione principale}: sensibilità specifica (vista);
  \item \textbf{origine}: retina dell'occhio;
  \item \textbf{destinazione}: diencefalo, tramite il chiasma ottico.
\end{itemize}

\rubrica{III nervo -- Oculomotore e IV nervo -- Trocleare}
\begin{itemize}
  \item \textbf{Funzione principale}: motoria, movimenti degli occhi;
  \item \textbf{origine}: mesencefalo;
  \item \textbf{destinazione}: il III innerva in modo somatico i muscoli
    retti superiore, inferiore e mediale, il muscolo obliquo inferiore e il
    muscolo elevatore della palpebra superiore; il IV innerva il muscolo
    obliquo superiore.
\end{itemize}

\rubrica{V nervo -- Trigemino}
\begin{itemize}
  \item \textbf{Funzione principale}: mista;
  \item \textbf{origine}: ramo oftalmico (sensitivo), ramo mascellare
    (sensitivo), ramo mandibolare (misto);
  \item \textbf{destinazione}: i tre rami raggiungono i nuclei sensitivi del
    ponte; il ramo mandibolare innerva inoltre i muscoli masticatori.
\end{itemize}

\rubrica{VI nervo -- Abducente}
\begin{itemize}
  \item \textbf{Funzione principale}: motoria, movimenti degli occhi;
  \item \textbf{origine}: ponte;
  \item \textbf{destinazione}: muscolo retto laterale.
\end{itemize}

\rubrica{VII nervo -- Faciale}
\begin{itemize}
  \item \textbf{Funzione principale}: mista;
  \item \textbf{origine}: sensitiva, dai recettori gustativi della lingua;
    motoria, dai nuclei motori del ponte;
  \item \textbf{destinazione}: fibre sensitive ai nuclei sensitivi del
    ponte; fibre motorie somatiche ai muscoli mimici.
\end{itemize}

\rubrica{VIII nervo -- Statoacustico}
\begin{itemize}
  \item \textbf{Funzione principale}: sensibilità specifica (equilibrio e
    udito);
  \item \textbf{origine}: recettori dell'orecchio interno (vestibolo e
    coclea);
  \item \textbf{destinazione}: nuclei vestibolare e cocleare del ponte e del
    bulbo.
\end{itemize}

\rubrica{IX nervo -- Glossofaringeo}
\begin{itemize}
  \item \textbf{Funzione principale}: mista;
  \item \textbf{origine}: sensitiva, dal terzo posteriore della lingua, da
    parte della faringe e del palato, e dalle arterie carotidi del collo;
    motoria, dai nuclei motori del bulbo;
  \item \textbf{destinazione}: fibre sensitive ai nuclei sensitivi del
    bulbo; fibre motorie somatiche ai muscoli faringei coinvolti nella
    deglutizione; fibre motorie viscerali alla ghiandola salivare parotide.
\end{itemize}

\rubrica{X nervo -- Vago}
\begin{itemize}
  \item \textbf{Funzione principale}: mista;
  \item \textbf{origine}: sensitiva viscerale da faringe, padiglione
    auricolare, canale uditivo esterno, diaframma e organi interni della
    cavità toracica e addomino-pelvica; motoria viscerale dai nuclei motori
    del bulbo;
  \item \textbf{destinazione}: fibre sensitive ai nuclei sensitivi e ai
    centri autonomi del bulbo; fibre motorie somatiche ai muscoli del
    palato e della faringe; fibre motorie viscerali agli organi degli
    apparati digerente, respiratorio e cardiovascolare, nella cavità
    toracica e addominale.
\end{itemize}

\rubrica{XI nervo -- Accessorio}
\begin{itemize}
  \item \textbf{Funzione principale}: motoria;
  \item \textbf{origine}: nuclei motori del midollo spinale e del bulbo;
  \item \textbf{destinazione}: il ramo interno innerva i muscoli volontari
    del palato, della faringe e della laringe; il ramo esterno controlla i
    muscoli sternocleidomastoideo e trapezio.
\end{itemize}

\rubrica{XII nervo -- Ipoglosso}
\begin{itemize}
  \item \textbf{Funzione principale}: motoria;
  \item \textbf{origine}: nuclei motori bulbari;
  \item \textbf{destinazione}: muscoli della lingua.
\end{itemize}
